\section{Індивідуальне завдання}

У межах індивідуального завдання розглядається cценарiй А. Backend-сервiс

\subsection{Основні вимоги}

Основні вимоги до криптографічного програмного забезпечення для цього сценарію
були визначені в попередньому розділі. Далі на їх основі обирається конкретна
криптографічна бібліотека та визначаються необхідні для Backend-сервісу
криптографічні механізми.

\subsection{Обрана криптографічна бібліотека}

Для реалізації криптографічних операцій обрано бібліотеку
\texttt{cryptography} для мови Python. Вона надає високорівневі інтерфейси
для типових криптографічних операцій, а також низькорівневий API для
спеціалізованих задач. Бібліотека підтримує симетричне та асиметричне
шифрування, цифрові підписи, геш-функції та функції виведення ключів\cite{cryptography}.

\subsection{Необхідні криптографічні функції}

Для захисту конфіденційних даних у Backend-сервісі доцільно використовувати
симетричне шифрування з автентифікацією, зокрема AES-GCM. Такий режим
забезпечує не лише конфіденційність, а й контроль цілісності та автентичності
зашифрованих даних.
Для операцій з відкритими ключами можуть використовуватися асиметричні
алгоритми, а для підтвердження автентичності повідомлень --- цифрові
підписи. Також необхідними є криптографічні геш-функції та функції виведення
ключів\cite{cryptography}.

\subsection{Робота з ключовою інформацією}

Оскільки сервер працює з конфіденційними даними, ключі не повинні
зберігатися безпосередньо у вихідному коді застосунку. Для їх зберігання
доцільно використовувати спеціалізовані засоби керування ключами та секретами.
Генерація ключів, nonce та інших випадкових значень повинна здійснюватися
за допомогою криптографічно стійкого генератора випадкових чисел. Крім
того, необхідно забезпечити контроль доступу до ключів та їх своєчасну
заміну відповідно до визначеного життєвого циклу \cite{OWASP}.

\subsection{Обґрунтування вибору бібліотеки}

Вибір бібліотеки \texttt{cryptography} зумовлений тим, що вона охоплює
основні криптографічні операції, необхідні для серверного застосунку, і
водночас надає високорівневий API для типових задач. Це дозволяє не
реалізовувати криптографічні алгоритми самостійно та зменшує ризик помилок
під час їх використання. Наявність підтримки AES-GCM, асиметричної криптографії, цифрових підписів
і функцій виведення ключів робить бібліотеку придатною для розглянутого
сценарію \cite{cryptography}.

\subsection{Ризики неправильного використання}

Основні ризики пов'язані не стільки з використанням самої бібліотеки,
скільки з неправильним застосуванням її можливостей. Наприклад, повторне
використання nonce з одним ключем під час роботи AES-GCM може призвести до
порушення безпеки шифрування \cite{cryptography}.
Також небезпечними є використання застарілих алгоритмів, слабких ключів,
неправильне зберігання ключової інформації та некоректне використання
низькорівневого API. Окремою проблемою є використання звичайних
криптографічних геш-функцій для зберігання паролів замість спеціалізованих
алгоритмів, призначених для цього завдання \cite{OWASP}.

